{{دیوان
| عنوان = فاطمہ بنت عبد اللہ
| ماخذ = ویکی ماخذ
| ماخذ_لنک = https://ur.wikisource.org/wiki/%D9%81%D8%A7%D8%B7%D9%85%DB%81_%D8%A8%D9%86%D8%AA_%D8%B9%D8%A8%D8%AF_%D8%A7%D9%84%D9%84%DB%81
| مقابلہ = اقبال سائبر لائبریری (اقبال اکادمی پاکستان)
| مقابلہ_لنک = https://www.iqbalcyberlibrary.net
}}
عرب لڑکی جو طرابلس کی جنگ میں غازیوں کو پانی پِلاتی ہوئی شہید ہُوئی

<poem>
فاطمہ! تُو آبرُوئے اُمّتِ مرحوم ہے
ذرّہ ذرّہ تیری مُشتِ خاک کا معصوم ہے

یہ سعادت، حُورِ صحرائی! تری قسمت میں تھی
غازیانِ دیں کی سقّائی تری قسمت میں تھی

یہ جہاد اللہ کے رستے میں بے تیغ و سِپَر
ہے جسارت آفریں شوقِ شہادت کس قدر

یہ کلی بھی اس گُلستانِ خزاں منظر میں تھی
ایسی چنگاری بھی یا رب، اپنی خاکستر میں تھی!

اپنے صحرا میں بہت آہُو ابھی پوشیدہ ہیں
بجلیاں برسے ہُوئے بادل میں بھی خوابیدہ ہیں!

فاطمہ! گو شبنم افشاں آنکھ تیرے غم میں ہے
نغمۂ عشرت بھی اپنے نالۂ ماتم میں ہے

رقص تیری خاک کا کتنا نشاط انگیز ہے
ذرّہ ذرّہ زندگی کے سوز سے لبریز ہے

ہے کوئی ہنگامہ تیری تُربتِ خاموش میں
پل رہی ہے ایک قومِ تازہ اس آغوش میں

بے خبر ہوں گرچہ اُن کی وسعتِ مقصد سے میں
آفرینش دیکھتا ہوں اُن کی اس مرقد سے میں

تازہ انجم کا فضائے آسماں میں ہے ظہور
دیدۂ انساں سے نامحرم ہے جن کی موجِ نور

جو ابھی اُبھرے ہیں ظُلمت خانۂ ایّام سے
جن کی ضَو ناآشنا ہے قیدِ صبح و شام سے

جن کی تابانی میں اندازِ کُہن بھی، نَو بھی ہے
اور تیرے کوکبِ تقدیر کا پرتَو بھی ہے
</poem>
